\documentclass[10pt]{article}

\usepackage{enumerate}
\usepackage[ruled,vlined,linesnumbered]{algorithm2e}


\begin{document}
	
\begin{center}
	\bf{Instituto de Computação - UFRJ
		
		Programação Orientada a Objetos (2025-1)
		
		12 de junho de 2025}
\end{center}
	
\bigskip

%%%%%%%%%%%%%%%%%%%%%%%%%%%%%%%%%%%%%%%%%%%%
\section{Trabalho 3 - Corrida}

\begin{itemize}
	
	\item Implementar a classe {\em Robot}.
	
	\item Considere um robô que encontra-se em um circuito de corrida de dimensões $lines \times columns$ tais que $lines$ e $columns$ são números positivos.
	
	\item Essa classe deve ter o método {\em move}, o qual, ao ser chamado, deve retornar UP, DOWN, LEFT, RIGHT ou STOP. Esse retorno significa o movimento que o robô deseja fazer para uma posição vizinha a atual na matriz que representa a sala.
		
	\item Inicialmente, o robô encontra-se na posição (1,1).
	
	\item Ele deve realizar 5 voltas no circuito com o menor número de movimentos.
	
	\item A linha de chegada é na posição $(\lfloor\frac{lines}{2}\rfloor,2)$.
	
	\item A localização atual do robô é representada por um par $(\ell', c')$ tal que $1 \le \ell' \le qttyLines$ e $1 \le c' \le qttyColumns$. O robô tem um GPS, para o qual ele pode perguntar qual a sua locali\-zação atual usando os métodos $getLine$ e $getColumn$. Esses métodos têm como parâmetro um inteiro que representa o {\em id} do robô.
	
\end{itemize}

\end{document}
